\chapter{Persistencia de sesión}
\label{ch:persistence}

\begin{note}
Stub. Las sesiones se serializan a archivos JSON bajo el
directorio de datos de app del SO. El formato es hoy
schema\_version $= 1$; un bump v2 para la migración MVCC P3
está planeado (Capítulo~\ref{ch:mvcc}).
\end{note}

\section{Layout en disco}

\begin{itemize}
    \item Windows: \code{\%APPDATA\%\textbackslash GraphHunter\textbackslash sessions\textbackslash <id>.json}
    \item Unix: \code{\$HOME/.local/share/GraphHunter/sessions/<id>.json}
\end{itemize}

\code{<id>} se valida contra \code{[A-Za-z0-9\_-]\{1,128\}}
para prevenir path traversal. El path se construye vía
\code{PathBuf::join}, nunca por concatenación de strings.

\section{Forma del SessionFile}

\begin{lstlisting}[style=rust]
pub struct SessionFile {
    pub id:              String,
    pub name:            String,
    pub created_at:      i64,
    pub entities:        Vec<Entity>,
    pub relations:       Vec<Relation>,
    pub path_node_ids:   Vec<String>,
    pub notes:           Vec<Note>,
    pub datasets:        Vec<DatasetInfo>,
    pub tags:            Vec<EntityTag>,
    // SIN ai_conversation -- el chat es deliberadamente efimero.
}
\end{lstlisting}

\section{Progreso de carga}

La carga es asíncrona. El \code{EventEmitter} recibe eventos
\code{SESSION\_LOAD\_PROGRESS} en cada etapa: parseo del archivo,
inserción de entidades, inserción de relaciones, reconstrucción
de índices (finalize del scorer), listo. La UI renderiza una
barra de progreso; HTTP / MCP simplemente ignoran los eventos
(\code{NoopEmitter}).

\section{Migración v2 (diferida)}

Cuando aterrice el cambio de writer MVCC P3, las aristas
heredadas (cargadas de archivos v1) deben backfillearse con un
\code{ingested\_at} sintético. El plan es usar el
\code{timestamp} de la arista como valor de backfill: no es
exactamente correcto (la arista fue ingerida en algún momento
posterior), pero es una elección definida y determinista que
permite a las queries de snapshot v2 producir resultados
sensatos contra datos v1. La UI muestra un banner
``snapshot aproximado'' para que el analista lo sepa.

\begin{inthecode}
\filepath{graph\_hunter\_api/src/state/persistence.rs} ---
load / save.\\
\filepath{graph\_hunter\_api/src/state/session\_types.rs} ---
\code{SessionFile}, \code{Note}, \code{EntityTag}.
\end{inthecode}
